\documentclass[11pt]{article}
\usepackage[utf8]{inputenc}
\usepackage{amsmath, amssymb, amsthm}
\usepackage[margin=1in]{geometry}

\theoremstyle{definition}
\newtheorem{definition}{Definition}
\newtheorem{theorem}{Theorem}
\newtheorem{lemma}{Lemma}

\theoremstyle{remark}
\newtheorem{remark}{Remark}

\title{Lecture 13: Rough SPDEs and the Kushner--Stratonovich Filtering Equation}
\date{08.01.2026}
\author{Tien Anh Vu}

\begin{document}
\maketitle

Good morning, everyone. Let us begin today's lecture. In this session, which is lecture thirteen on Stochastic Partial Differential Equations (SPDEs), we will analyze a specific non-linear SPDE and discuss the necessity of rough path techniques to make sense of the integrals that arise from it.

\section*{1. Main Equation}
Let us start by writing down our main equation for today. We consider the following stochastic partial differential equation on the domain $L^2([0,2\pi])$:
\[
du=\partial_{xx}^{2}u\,dt+f(u)\,dt+g(u)\partial_xu\,dt+\sigma\,dW_t.
\]

To understand this equation, let us break it down term by term. The left-hand side, $du$, represents the stochastic differential of our solution process $u(t,x)$. On the right-hand side, the first term $\partial_{xx}^{2}u\,dt$ is the standard second-order spatial derivative, corresponding to the heat operator or Laplacian in one dimension. The second term, $f(u)\,dt$, is a non-linear reaction term. The third term, $g(u)\partial_xu\,dt$, represents a non-linear transport or gradient component. Finally, we have the additive noise term $\sigma\,dW_t$, where $W_t$ denotes the driving Wiener process.

\section*{2. Semigroup and Decomposition}
To approach this equation, we introduce the heat semigroup, denoted by $P_t$, which is generated by the spatial second derivative:
\[
P_t=e^{t\partial_{xx}^{2}}.
\]

Instead of tackling $u(t,x)$ directly, it is highly beneficial to decompose our solution into a smoother component and a rougher stochastic component. We write the solution as:
\[
u(t,x)=v(t,x)+\sigma W_{\partial_{xx}^{2}}(t).
\]
Here, the term $W_{\partial_{xx}^{2}}(t)$ represents the stochastic convolution. It is defined as the integral $\int_0^t e^{(t-s)\partial_{xx}^{2}}\,dW_s$. Due to the properties of the white noise and the heat semigroup, this stochastic convolution lives in the fractional Sobolev space $H^{\alpha,2}$ for regularity parameter $\alpha<\frac{1}{2}$.

By substituting this decomposition back into our original SPDE and applying the variation of constants formula, we can write down an integral equation for the smoother part, $v(t,x)$:
\[
v(t,x)=P_tu_0(x)+\int_0^t P_{t-s}\bigl(f(u(s,\cdot))+g(u(s,\cdot))\partial_xv(s,\cdot)\bigr)\,ds+\int_0^t P_{t-s}\bigl(g(u(s,\cdot))\partial_xW_{\partial_{xx}^{2}}(s,\cdot)\bigr)(x)\,ds.
\]

Let us read through this equation carefully. The first term, $P_tu_0(x)$, is simply the initial condition $u_0(x)$ propagated forward in time by the heat semigroup. The second term is the time integral from $0$ to $t$ of the semigroup applied to our smooth non-linearities, incorporating both the reaction term $f(u)$ and the gradient of our smooth part $\partial_xv$.

\section*{3. The Problematic Spatial Integral}
The critical part of this equation is the final term. Let us expand the spatial action of the semigroup in this last integral over our domain $[0,2\pi]$:
\[
\int_0^{2\pi}P_{t-s}(x-y)g(u(s,y))\partial_yW_{\partial_{xx}^{2}}(s,y)\,dy.
\]
Because the stochastic convolution has low spatial regularity $(\alpha<\frac{1}{2})$, taking its spatial derivative $\partial_yW_{\partial_{xx}^{2}}$ results in a distribution rather than a classical function. Consequently, this integral cannot be understood in the standard Lebesgue or Riemann sense. It is, fundamentally, a rough integral with respect to the spatial variable $y$.

To rigorously define such an object, we must take an interlude and briefly discuss the principles of rough path theory.

\section*{4. Rough Paths Interlude}
Let us consider two real-valued processes $X$ and $Y$ defined on a time interval $[0,T]$. Assume that both $X$ and $Y$ belong to the H\"older space $C^\alpha$, where the H\"older exponent $\alpha$ is strictly between $\frac{1}{3}$ and $\frac{1}{2}$.

We want to define the integral of $Y$ with respect to $X$:
\[
\int_s^t Y_r\,dX_r
:=
\lim_{|\mathcal P|\to 0}\sum_{i=0}^{n-1}Y_{t_i}\bigl(X_{t_{i+1}}-X_{t_i}\bigr).
\]
Here, we are attempting to take the limit of Riemann-Stieltjes sums over partitions $\mathcal P$ of the interval $[s,t]$ as the mesh size of the partition $\mathcal P$ goes to zero. However, suppose $Y_t$ is a function of $X_t$, say $Y_t=f(X_t)$ for some continuously differentiable function $f\in C^1(\mathbb R)$. Because $\alpha<\frac{1}{2}$, the sum of the H\"older regularities of the integrand and integrator is strictly less than $1$. Standard Young integration theory tells us that this specific partition-limit expression for $\int_s^t Y_r\,dX_r$ need not exist.

To overcome this lack of convergence, we expand the integrand further. Here the first-order approximation is $Y_s(X_t-X_s)$, obtained by freezing the integrand at the left endpoint $s$. Let us look at the difference between the integral and this first-order approximation:
\[
\int_s^t Y_r\,dX_r-Y_s(X_t-X_s)=\int_s^t (Y_r-Y_s)\,dX_r.
\]
By expressing the increment $(Y_r-Y_s)$ as an integral of its derivative, $\int_s^r f'(X_u)\,dX_u$, we can substitute this back into our equation:
\[
=\int_s^t\int_s^r f'(X_u)\,dX_u\,dX_r.
\]
We can now freeze the inner derivative at the lower limit $s$:
\[
\int_s^r f'(X_u)\,dX_u
=
f'(X_s)(X_r-X_s)+\int_s^r\bigl(f'(X_u)-f'(X_s)\bigr)\,dX_u.
\]
Substituting this decomposition gives
\[
\int_s^t\int_s^r f'(X_u)\,dX_u\,dX_r
=
\underbrace{f'(X_s)\int_s^t (X_r-X_s)\,dX_r}_{\text{leading rough term}}
+\underbrace{\widetilde R_{s,t}}_{\text{higher-order remainder}},
\]
where
\[
\widetilde R_{s,t}
:=
\int_s^t\int_s^r\bigl(f'(X_u)-f'(X_s)\bigr)\,dX_u\,dX_r.
\]
So $\widetilde R_{s,t}$ is the higher-order remainder generated by replacing $f'(X_u)$ with $f'(X_s)$ in the inner integral.

With our notation,
\[
\mathbb{X}_{s,t}:=X_{s,t}^{(2)}=\int_s^t (X_r-X_s)\,dX_r,
\qquad
\mathbb{Z}_{s,t}:=f'(X_s)\,\mathbb{X}_{s,t}.
\]
So $\mathbb{X}_{s,t}$ is the second level of the rough path, while $\mathbb{Z}_{s,t}$ is the leading term in the expansion.
Quantitatively, the remainder is of order $|t-s|^{3\alpha}$ ($3\alpha$-H\"older regularity), compared to the leading term order $|t-s|^{2\alpha}$ ($2\alpha$-H\"older regularity), so it has higher regularity.
The central insight of rough path theory is that this leading rough term cannot be handled by classical limits when paths are too irregular; one needs the rough-path second-level information to make it rigorous.

\section*{5. Back to the SPDE}
We will apply these precise rough path concepts to handle the problematic spatial integral in our SPDE. We will continue from this point in the next section.

\subsection*{5.1. Convergence Mechanism for the Rough Integral}
Let us continue from where we left off. We established that to make sense of integrals where both the integrand and the integrator have low regularity, specifically when the H\"older exponent $\alpha$ is between $\frac13$ and $\frac12$, we must introduce an additional object. We denoted this as the iterated integral $\mathbb{X}_{s,t}$.

We can achieve convergence of Riemann--Stieltjes type sums over a partition $\mathcal P$ of the interval. Specifically, we look at the limit as the mesh size $|\mathcal P|$ goes to zero of
\[
\lim_{|\mathcal P|\to0}
\sum_{[t_i,t_{i+1}]\in\Pi(\mathcal P)}
\Bigl(
f(X_{t_i})(X_{t_{i+1}}-X_{t_i})
+f'(X_{t_i})\mathbb{X}_{t_i,t_{i+1}}
\Bigr),
\]
where $\Pi(\mathcal P)$ denotes the family of subintervals of the partition.
Equivalently, this defines the rough integral by
\[
\int_s^t f(X_r)\,dX_r
:=
\lim_{|\mathcal P|\to0}
\sum_{[t_i,t_{i+1}]\in\Pi(\mathcal P)}
\Bigl(
f(X_{t_i})(X_{t_{i+1}}-X_{t_i})
+f'(X_{t_i})\mathbb{X}_{t_i,t_{i+1}}
\Bigr).
\]

For this limit to converge and provide a robust definition of the rough integral, the a priori defined object $\mathbb{X}_{s,t}$ must satisfy two critical conditions. First, it must have $2\alpha$-H\"older regularity. Second, it must satisfy the consistency condition known as Chen's relation. For intermediate times $u$ with $s\le u\le t$,
\[
\mathbb{X}_{s,t}-\mathbb{X}_{s,u}-\mathbb{X}_{u,t}
=(X_u-X_s)(X_t-X_u).
\]
It is called a consistency condition because it enforces compatibility of increments on subintervals with the increment on the whole interval $[s,t]$.
Meaning: the second-level increment on $[s,t]$ is consistent with splitting at $u$. It is the key algebraic condition that makes the rough integral independent of partition refinement.

\subsection*{5.2. Remainder Estimate}
Assuming our function $f$ is twice continuously differentiable, $f\in C^2$, we define the remainder $\widetilde R_{s,t}$ as the difference between the exact integral and its second-order expansion:
\[
\widetilde R_{s,t}
:=
\int_s^t f(X_r)\,dX_r
-f(X_s)(X_t-X_s)
-f'(X_s)\mathbb{X}_{s,t}.
\]
Equivalently,
\[
\int_s^t f(X_r)\,dX_r
=
f(X_s)(X_t-X_s)+f'(X_s)\mathbb{X}_{s,t}+\widetilde R_{s,t}.
\]
According to Lemma A.13 in the appendix, one has
\[
|\widetilde R_{s,t}|
\le
K\|f\|_{C_b^2}
\Bigl(
\|X\|_{C^\alpha}^3
+\|X\|_{C^\alpha}\|\mathbb{X}\|_{C^{(2\alpha)}}
\Bigr)
|t-s|^{3\alpha}.
\]
The crucial factor is $|t-s|^{3\alpha}$.

Because $\alpha>\frac13$, we have $3\alpha>1$. Write $3\alpha=1+\delta$ with $\delta>0$. Summing over partition intervals,
\[
\sum_{[t_i,t_{i+1}]\in\Pi(\mathcal P)} |t_{i+1}-t_i|^{1+\delta}
\le
|\mathcal P|^\delta\,\mathrm{Var}(\mathrm{Id},[0,T]).
\]
Since $\mathrm{Var}(\mathrm{Id},[0,T])=T$, the right-hand side tends to $0$ as $|\mathcal P|\to0$. This vanishing remainder is the analytic mechanism ensuring convergence of the rough integral.

\subsection*{5.3. Concrete Form of the Defining Term}
To ground this structure, consider the defining term
\[
\int_s^t (X_r-X_s)\,dX_r.
\]
If $X$ is differentiable, $dX_r=\dot X_r\,dr$, then
\[
\int_s^t (X_r-X_s)\dot X_r\,dr
=
\frac12(X_t^2-X_s^2)-X_s(X_t-X_s)
=
\frac12(X_t-X_s)^2.
\]

If $X$ is standard Brownian motion and the integral is interpreted in the It\^o sense,
\[
\int_s^t (X_r-X_s)\,dX_r
=
\int_s^t X_r\,dX_r
-X_s(X_t-X_s).
\]
Using It\^o's formula,
\[
\int_s^t X_r\,dX_r
=
\frac12(X_t^2-X_s^2)-\frac12(t-s),
\]
hence
\[
\int_s^t (X_r-X_s)\,dX_r
=
\frac12(X_t-X_s)^2-\frac12(t-s).
\]
The extra term $-\frac12(t-s)$ is the quadratic variation correction. This illustrates that the algebraic object $\mathbb{X}_{s,t}$ is not determined by $X$ alone unless one specifies the integration convention (It\^o vs Stratonovich).

\subsection*{5.4. Consequence for the SPDE}
By defining $\mathbb{X}_{s,t}$ appropriately in the spatial domain for the stochastic convolution, we can rigorously define the nonlinear transport structure corresponding to $g(u)\partial_xu$ in the SPDE discussed in the previous segment, thereby completing the formulation.

\section*{6. Stationary Solution for a Linear SPDE}
Let us proceed. Now that we have established the necessary rough path machinery for dealing with non-linear terms and low-regularity spatial noise, I want us to shift our focus slightly. We will look at how we can analyze the long-time behavior of these systems by studying stationary solutions, focusing on a specific linear SPDE known as the stochastic Ornstein--Uhlenbeck equation.

Let us write down the governing equation for this process:
\[
dX_t=(\partial_{xx}-1)X_t\,dt+\sigma\,dW_t.
\]
Notice the modification here compared to the standard heat equation we discussed previously. We have introduced a $-1$ into the linear differential operator, resulting in the operator $(\partial_{xx}-1)$. This addition is crucial because it acts as a damping term, which gives rise to exponential decay in our system.

Consequently, the semigroup associated with this modified operator, which we will denote here as $P_t$, contains this exponential decay factor:
\[
P_t=e^{(\partial_{xx}-1)t}=e^{-t}e^{t\partial_{xx}}.
\]

To construct a stationary solution, a solution whose statistical properties do not change over time, we cannot simply start the process at time $t=0$. Instead, we must imagine that the process has been running since the infinite past. We formulate the stochastic convolution as an integral from $-\infty$ to $t$:
\[
X_t
=
\sigma\int_{-\infty}^t e^{(t-s)(\partial_{xx}-1)}\,dW_s
=
\sigma\int_{-\infty}^t e^{-(t-s)}P_{t-s}\,dW_s.
\]

To make sense of an integral starting from $-\infty$, we must utilize a two-sided Wiener process. We define this process, $W_t$, for all real numbers $t$, typically by pasting together two independent standard Wiener processes at the origin:
\[
W_t=
\begin{cases}
W_t^f, & t\ge0,\\
W_{-t}^b, & t<0.
\end{cases}
\]
With this two-sided process, the increments are well-defined. For the mixed case $s<0<t$,
\[
W_t-W_s
=
\bigl(W_t^f-W_0^f\bigr)+\bigl(W_0^b-W_{-s}^b\bigr)
\sim
\mathcal N\!\bigl(0,(t-s)\mathrm{Id}\bigr).
\]

To see why this integral from $-\infty$ represents the stationary state, let us contrast it with the standard initial value problem starting at $t=0$. If we solve the equation forward from a given initial state $X_0$, the variation of constants formula gives us:
\[
X_t
=
e^{-t}P_tX_0
+\sigma\int_0^t e^{-(t-s)}P_{t-s}\,dW_s.
\]
Because of the $e^{-t}$ term, the influence of the initial condition $X_0$ decays exponentially to zero as $t\to\infty$. What remains is the stochastic integral, which converges in distribution to our stationary solution.

We can explicitly characterize this stationary distribution. Because the equation is linear and driven by additive Gaussian noise, the stationary solution $X_t$ is a Gaussian random field with mean zero. Its spatial covariance is determined by the integral of the squared semigroup operator over all time. We can define this stationary distribution as:
\[
X_t\sim N\!\left(0,\frac{\sigma^2}{2}(I-\partial_{xx})^{-1}\right).
\]
The operator $\frac{\sigma^2}{2}(I-\partial_{xx})^{-1}$ is the covariance operator of our stationary field. In harmonic analysis, we have the result that the action of this inverse operator can be represented as an integral operator with a spatial kernel $K$:
\[
\frac{\sigma^2}{2}(I-\partial_{xx})^{-1}f(x)
=
\int_0^{2\pi} f(y)K(x-y)\,dy.
\]

To find the explicit form of this kernel $K(x)$, we analyze how the semigroup acts on the Fourier basis functions. The action of the standard heat semigroup $P_{2t}$ on a cosine function of wave number $k$ is given by:
\[
\bigl(P_{2t}\cos(k\cdot)\bigr)(x)=e^{-2k^2t}\cos(kx).
\]
By integrating these decaying exponential modes over time, one can derive the Fourier series representation for the kernel $K(x)$. We find that $K(x)$ is given by the following infinite sum over the wave numbers $k$:
\[
K(x)=\frac{\sigma^2}{2\pi}\sum_{k\in\mathbb Z}\frac{\cos(kx)}{1+k^2}.
\]
This series can be evaluated explicitly. The closed-form expression for the spatial covariance kernel of our stationary solution is:
\[
K(x)=\frac{\sigma^2}{2}\,\frac{\cosh(|x|-\pi)}{\sinh(\pi)}.
\]
This explicit kernel tells us precisely how the random fluctuations at a point $x$ are correlated with the fluctuations at a point $y$ in the stationary regime of the stochastic Ornstein--Uhlenbeck process. The exponential nature of the hyperbolic functions reflects how the correlation decays with spatial distance across our domain $[0,2\pi]$.


\section*{7. Spatial Regularity from the Stationary Covariance Kernel}
\subsection*{7.1. Setup and Kolmogorov Criterion}
Let us resume our discussion. In the previous segment, we derived the explicit stationary covariance kernel $K(x)$ for the stochastic Ornstein--Uhlenbeck process. Now, we will utilize this kernel to deduce the spatial regularity of the paths of our stationary solution $X_t(x)$.

Let us define our spatial domain precisely. We will consider $x\in[-\pi,\pi]$ with periodic extension, which functions identically to our earlier interval of $[0,2\pi]$.

To rigorously establish the continuity and regularity of the paths $x\mapsto X_t(x)$, we will apply Kolmogorov's continuity theorem. The theorem states that if there exist constants $p\ge1$, $\varepsilon>0$, and $K>0$ such that
\[
\mathbb E\bigl(|X_t(x)-X_t(y)|^p\bigr)\le K|x-y|^{1+\varepsilon},
\]
then the process admits a continuous modification, and its paths are almost surely locally H\"older continuous. Specifically, for an appropriate H\"older exponent $\alpha$, we have
\[
\sup_{x,y}\frac{|X_t(x)-X_t(y)|}{|x-y|^\alpha}<\infty
\qquad\text{almost surely}.
\]

To apply this theorem, we examine spatial increments of the stationary solution. We know that $X_t$ is a mean-zero Gaussian random field with covariance operator $C$, i.e.
\[
X_t\sim N(0,C).
\]

\subsection*{7.2. Increment Moment via Covariance Operator}
Start with the second moment of an increment. In weak form, with test functions $f,g$:
\[
\mathbb E\bigl(\langle X_t,f-g\rangle^2\bigr).
\]
Because for centered Gaussian variables the second moment is the variance,
\[
\mathbb E\bigl(\langle X_t,f-g\rangle^2\bigr)
=
\int_0^{2\pi}(f-g)(x)\,C(f-g)(x)\,dx.
\]
Using the covariance kernel representation of $C$:
\[
=
\int_0^{2\pi}\!\!\int_0^{2\pi}
(f-g)(\tilde x)\,K(\tilde x-\tilde y)\,(f-g)(\tilde y)\,d\tilde y\,d\tilde x.
\]

\subsection*{7.3. Dirac Evaluation and Exact Increment Formula}
To recover point evaluations, take $f\to\delta_x$ and $g\to\delta_y$. Then
\[
\int_0^{2\pi}X_t(z)\,\delta_x(dz)\to X_t(x),
\qquad
\int_0^{2\pi}X_t(z)\,\delta_y(dz)\to X_t(y).
\]
Hence
\[
\mathbb E\bigl(|X_t(x)-X_t(y)|^2\bigr)
=
\int_0^{2\pi}\!\!\int_0^{2\pi}
(\delta_x-\delta_y)(\tilde x)\,K(\tilde x-\tilde y)\,(\delta_x-\delta_y)(\tilde y)\,d\tilde y\,d\tilde x.
\]

Evaluate first in $\tilde y$:
\[
=
\int_0^{2\pi}
(\delta_x-\delta_y)(\tilde x)\,[K(\tilde x-x)-K(\tilde x-y)]\,d\tilde x.
\]
Then in $\tilde x$:
\[
=[K(x-x)-K(x-y)]-[K(y-x)-K(y-y)].
\]
Using $K(0)$ and symmetry $K(y-x)=K(x-y)$:
\[
=2K(0)-2K(x-y).
\]
Therefore,
\[
\mathbb E\bigl(|X_t(x)-X_t(y)|^2\bigr)
=
2\bigl(K(0)-K(x-y)\bigr).
\]

\subsection*{7.4. Kolmogorov Conclusion}
To connect with Kolmogorov's criterion, inspect $K(0)-K(x-y)$. From the explicit kernel, this difference is Lipschitz, so
\[
\mathbb E\bigl(|X_t(x)-X_t(y)|^2\bigr)\le L|x-y|.
\]
For $p=2$ this gives exponent $1$, while Kolmogorov requires strictly larger than $1$. Since increments are Gaussian, higher even moments are controlled by powers of the second moment; e.g. for $p=4$:
\[
\mathbb E\bigl(|X_t(x)-X_t(y)|^4\bigr)\lesssim |x-y|^2,
\]
which satisfies the $1+\varepsilon$ condition. Hence Kolmogorov's theorem yields almost sure spatial H\"older continuity of stationary paths.




\section*{8. Rough Lift and Well-Posedness Framework}
\subsection*{8.1. Rough Lift of the Spatial Process}
Let us continue. Having established the spatial regularity of the stochastic Ornstein--Uhlenbeck process in the previous segment, we can now definitively state that its paths possess H\"older regularity strictly less than one-half. This allows us to apply the rough path framework directly to the spatial dimension of our stochastic partial differential equation.

Consequently, for any fixed time $t$, the spatial process $X_t$ can be lifted to a rough path over our domain. We denote this pair consisting of the path and its iterated integral as
\[
\bigl(X_t(x),\mathbb X_t(x,y)\bigr),
\qquad 0\le x<y\le 2\pi.
\]
To rigorously construct this iterated integral $\mathbb X_t(x,y)$, we must employ an approximation procedure, as the direct Riemann--Stieltjes integral is ill-defined. We define it as the limit as $\varepsilon\downarrow0$:
\[
\mathbb X_t(x,y)
:=
\lim_{\varepsilon\downarrow0}
\int_x^y\bigl(X_t^{\varepsilon}(v)-X_t^{\varepsilon}(x)\bigr)\,dX_t^{\varepsilon}(v).
\]
Here, we integrate over an approximating sequence of Gaussian processes $(X_t^{\varepsilon}(x))_{\varepsilon}$, smooth in the spatial variable $x$. For this limit to be valid, the sequence is chosen so that the approximating processes converge to the true process in the mean-square sense.

The notation $(X_t^{\varepsilon}(x))_{\varepsilon}$ means a family of approximating processes labeled by the smoothing parameter $\varepsilon>0$. We index by $\varepsilon$ because each $\varepsilon$ gives a different smooth approximation of the same rough object $X_t(x)$, and then we take $\varepsilon\downarrow 0$ to recover the original rough process. Thus $\varepsilon$ is not time; it is the approximation scale (regularization level).

\subsection*{8.2. Temporal Continuity of the Lift}
Next, we consider continuity of these lifted paths with respect to time. For every spatial H\"older exponent $\alpha<\frac12$, there exists a strictly positive constant $\theta>0$ such that
\[
\mathbb E\!\left(
\|X_t-X_s\|_{C^{\alpha}}^{q}
+
\|\mathbb X_t-\mathbb X_s\|_{C^{(2\alpha)}}^{q/2}
\right)
\le C_q\,|t-s|^{\theta q},
\qquad \forall q\ge1.
\]
In words, this says the lifted process is time-H\"older in moment: when $t$ is close to $s$, both the first level $X_t$ and second level $\mathbb X_t$ are close to their values at time $s$, with expected increment size controlled by a positive power of $|t-s|$.
To be absolutely clear on the notation used in this bound, the H\"older seminorm of the first-order path is
\[
\|u\|_{\alpha}
:=
\sup_{x\ne y}\frac{|u(x)-u(y)|}{|x-y|^{\alpha}},
\]
and the seminorm for the second-order iterated integral is
\[
\|U\|_{2\alpha}
:=
\sup_{x<y}\frac{|U(x,y)|}{|x-y|^{2\alpha}}.
\]

\subsection*{8.3. Local and Global Well-Posedness}
With this well-defined rough path structure in place, we state the central result.
\begin{theorem}[Local and global well-posedness]
Consider the SPDE
\[
du=\bigl(\partial_{xx}u+f(u)+g(u)\partial_xu\bigr)\,dt+\sigma\,dW_t.
\]
with initial condition $u(0,\cdot)=u_0$.
Assume $\beta\in(\frac13,\frac12)$ and $u_0\in C^{\beta}$. Then, for almost every realization of the rough path, there exists a strictly positive terminal time $T>0$ such that there is a unique mild solution in
\[
C([0,T];C^{\beta}).
\]
Furthermore, if $g$ is bounded and all derivatives of $f$ and $g$ are bounded, then this local solution is global and independent of the specific rough path realization.
\end{theorem}

\subsection*{8.4. Picard Iteration Structure}
The standard proof idea relies on Picard iteration, adapted to the rough path topology outlined above. To execute this Picard iteration, we return to the variation-of-constants formulation for the smoother component $v(t,x)$ that isolates the rough noise:
\[
v(t,x)
=
P_tu_0(x)
+
\int_0^t P_{t-s}\bigl(f(u(s,\cdot))+g(u(s,\cdot))\partial_xv(s,\cdot)\bigr)(x)\,ds.
\]
By controlling the smooth part $v(t,x)$ via classical semigroup estimates and treating rough integration against the stochastic convolution through the bounds on $\mathbb X_t(x,y)$, the Picard map is a contraction, yielding the theorem.

\section*{9. SPDE Bounds and Closing Estimates}
Let us carefully review and formalize the bounds for our stochastic partial differential equation, and then write out the precise formulation of the filtering problem. We must ensure every equation and derivation step is rigorously documented.

\subsection*{9.1. Decomposition and Step-1 Bound}
We will begin by concluding our estimates for the SPDE. Recall that we utilize a decomposition where the solution $u$ is separated into a smoother component $v$ and the stochastic convolution $W_{\partial_{xx}^{2}}$, written as
\[
u=v+W_{\partial_{xx}^{2}}.
\]
To control the solution iteratively, we define Step 1 as the analysis of an ordinary differential equation with a Lipschitz continuous coefficient. We denote this classical, smooth part of the approximation operator as $M^{(1)}_{T,X}$.

We establish a supremum bound over time $t\le T$ for the $C^1$ norm of this operator applied to our smooth component $V(t)$. The estimate is bounded by a constant $C_K$ multiplied by the rough path norm:
\[
\|M^{(1)}_{T,X}\|_{1,T}
=
\sup_{t\le T}\|M^{(1)}_{T,X}V(t)\|_{C^1}
\le
C_K(1+\|\mathbf X\|).
\]
This bound is valid under the assumption that the $C^\beta$ norm of our initial condition $u_0$, the $C^1$ norm of our non-linear function $f$, and the supremum norm of $V$ over time $T$ are all bounded by a constant $K$:
\[
\|u_0\|_{C^\beta}\le K,\qquad
\|f\|_{C^1}\le K,\qquad
\|V\|_{T,\infty}\le K.
\]
It is necessary to define the rough path norm $\|\mathbf X\|$ explicitly. For a fixed spatial regularity parameter $\alpha\in(1/3,1/2)$, the norm over the unit time interval is defined as the supremum of the sum of the H\"older norms of the first-order path $X_t$ and the second-order iterated integral $\mathbb X_t$:
\[
\|\mathbf X\|
=
\sup_{t\le1}\bigl(\|X_t\|_{C^\alpha}+\|\mathbb X_t\|_{C^{2\alpha}}\bigr).
\]

\subsection*{9.2. Heat-Semigroup Estimates}
To properly close our estimates in the integral equation, we need to control the spatial derivative of the heat semigroup applied to our functions, $\partial_xP_tf$, and obtain what we call a ``good estimate''.
First, we rely on the continuity property of the heat kernel, which preserves the H\"older regularity of a function. We have the bound:
\[
\|P_tf\|_{C^\beta}\le\|f\|_{C^\beta}.
\]
Second, for any function $f$ in the periodic H\"older space $C^\beta_{\mathrm{per}}$, the heat semigroup smooths the function, allowing us to bound its spatial gradient. For all time $t>0$, the supremum norm is bounded by:
\[
\|\nabla P_tf\|_{\infty}
\le
C_t\,t^{2\beta-1}\|f\|_{C^\beta}.
\]
The integrable singularity $t^{2\beta-1}$ is the critical component that allows these integrals to converge over time.

\section*{10. Filtering Problem Formulation}
Now, let us transition to the Filtering Problem.

\subsection*{10.1. Signal and Observation Models}
We define a hidden state process, which we refer to as the signal. This process evolves according to a stochastic differential equation, labeled (S):
\[
\text{(S)}:\quad dX_t=b(X_t)\,dt+\sigma(X_t)\,dW_t.
\]
We specify that the initial state $X_0$ is distributed according to an initial probability measure $P_0$, written as
\[
X_0\sim P_0.
\]
Because $X_t$ is hidden from us, we must rely on an observation process, labeled (O), which is governed by its own differential equation:
\[
\text{(O)}:\quad dY_t=h(X_t)\,dt+\rho\,dV_t.
\]

\subsection*{10.2. Filter Definition}
Our objective is to determine the conditional probability distribution of the state $X_t$, given the continuous history of observations up to time $t$, which we denote as the sigma-algebra $\mathcal Y_{0,t}$. We define this filter measure as $\eta_t^Y$:
\[
\eta_t^Y=\mathbb P(X_t\mid \mathcal Y_{0,t}).
\]
To evaluate a test function $f$ against our current estimate of the state, we compute the conditional expectation by integrating $f$ with respect to our filter measure $\eta_t^Y$:
\[
\mathbb E\!\bigl(f(X_t)\mid \mathcal Y_{0,t}\bigr)=\int f\,d\eta_t^Y.
\]

\subsection*{10.3. Kushner--Stratonovich Equation and Innovation}
The temporal evolution of this conditional expectation is given by the Kushner--Stratonovich equation. We write the stochastic differential for the filter applied to the test function $f$ as follows:
\[
\underbrace{d\eta_t^Y(f)}_{\text{update of filtered expectation}}
=
\underbrace{\eta_t^Y(Lf)\,dt}_{\text{prediction / internal dynamics}}
+
\underbrace{\operatorname{Cov}_{\eta_t^Y}(f,h)(\rho\rho^T)^{-1}\bigl(dY_t-\eta_t^Y(h)\,dt\bigr)}_{\text{measurement correction via innovation}}.
\]
Let us examine the two primary components of this equation. The first term on the right-hand side, $\eta_t^Y(Lf)\,dt$, represents the weak form of the Fokker--Planck Equation (FPE). It dictates how our belief of the state distribution would evolve according to the signal's internal dynamics, represented by the generator $L$, in the absence of new observations.
The second term represents the measurement update. The differential term $\bigl(dY_t-\eta_t^Y(h)\,dt\bigr)$ is defined as the innovation process. It quantifies the difference between the actual observation increment $dY_t$ and our expected observation increment $\eta_t^Y(h)\,dt$. This new information is then scaled by the inverse of the observation noise covariance, $(\rho\rho^T)^{-1}$, and the conditional covariance between our test function $f$ and the observation function $h$, $\operatorname{Cov}_{\eta_t^Y}(f,h)$, to continuously correct our state estimate.

\subsection*{10.4. Density (Strong) Form of Kushner--Stratonovich}
Now, we will take a step further. We assume that our filter measure, $\eta_t^Y$, admits a density with respect to the Lebesgue measure. For simplicity of notation, we will also denote this conditional probability density as $\eta_t^Y(x)$. Our objective is to derive the strong form of the Kushner-Stratonovich equation, which is a stochastic partial differential equation that governs the direct spatial and temporal evolution of this density profile.

\subsubsection*{10.4.1. Formal and Differential Forms}
Let us write down this equation for the density, as written here in the notes:
\[
\underbrace{\partial_t\eta_t^Y(x)}_{\text{time evolution of filter density}}
=
\underbrace{L^\ast\eta_t^Y(x)}_{\text{prediction / Fokker--Planck term}}
+
\underbrace{\bigl(h(x)-\eta_t^Y(h)\bigr)\eta_t^Y(x)(\rho\rho^T)^{-1}dY_t}_{\text{observation-driven correction (formal form)}}.
\]
While this is written using a formal partial time derivative $\partial_t$ for brevity, in rigorous stochastic calculus, we express this as a stochastic differential equation driven by the innovation process. Let us write the rigorously expanded differential form:
\[
\underbrace{d\eta_t^Y(x)}_{\text{stochastic update of density}}
=
\underbrace{L^\ast\eta_t^Y(x)\,dt}_{\text{prediction drift}}
+
\underbrace{\bigl(h(x)-\eta_t^Y(h)\bigr)\eta_t^Y(x)(\rho\rho^T)^{-1}\bigl(dY_t-\eta_t^Y(h)\,dt\bigr)}_{\text{innovation-based measurement update}}.
\]

\subsubsection*{10.4.2. Term-by-Term Interpretation}
Let us break down this equation carefully, term by term, to understand the mechanics of the filter density.

The left-hand side, $d\eta_t^Y(x)$ (or formally $\partial_t\eta_t^Y(x)$), represents the stochastic update to the probability density of our hidden state $x$ at time $t$, conditioned on the continuous stream of observations up to that time.

The first term on the right-hand side is $L^\ast\eta_t^Y(x)$. Notice that we are no longer using the generator $L$ as we did in the weak form. Instead, we use its formal adjoint, $L^\ast$. This operator is often referred to as the Fokker-Planck operator or the forward Kolmogorov operator. The mathematical reason for this shift is integration by parts: when we move from the weak formulation, where the generator $L$ acts on a smooth test function $f$, to the strong formulation for the density, we transfer the spatial derivatives from the test function onto the density itself, transforming $L$ into $L^\ast$. Physically, this term describes the internal dynamics--the natural diffusion and drift of the state's probability density--as it would evolve between observations.

The second part of the equation is the measurement update. We have the algebraic expression $\bigl(h(x)-\eta_t^Y(h)\bigr)\eta_t^Y(x)$. This pointwise expression replaces the abstract conditional covariance $\operatorname{Cov}_{\eta_t^Y}(f,h)$ from our weak formulation. It scales the current density at state $x$, $\eta_t^Y(x)$, by the difference between the observation function evaluated exactly at that state, $h(x)$, and the globally expected observation, $\eta_t^Y(h)$.

This update vector is then weighted by the inverse of the observation noise covariance matrix, $(\rho\rho^T)^{-1}$. This matrix inverse serves as a precision weight; it quantifies our confidence in the incoming measurements. If the observation noise $\rho$ is very small, this inverse matrix becomes large, meaning our density will react aggressively to correct itself based on new data.

Finally, this entire corrective term is multiplied by the innovation increment, $\bigl(dY_t-\eta_t^Y(h)\,dt\bigr)$ (formally represented by the observation differential in the notes). As discussed, this innovation is the mathematical representation of the ``surprise'' or the strictly new information provided by the observation $dY_t$ that was not already predicted by our current estimate $\eta_t^Y(h)\,dt$.

\subsubsection*{10.4.3. Nonlinearity and Nonlocal Dependence}
It is vital to recognize a fundamental mathematical challenge presented by this SPDE: it is a highly non-linear equation. The non-linearity does not come from the state dynamics, but from the filtering process itself. The term $\eta_t^Y(h)$, which appears twice in our measurement update, is the expected value of the observation function. It is computed by integrating over the entire state space:
\[
\eta_t^Y(h)=\int h(z)\eta_t^Y(z)\,dz.
\]

Because this integral depends on the entire density profile $\eta_t^Y$, the evolution of the density at any specific local point $x$ depends non-locally on the density everywhere else. This integral non-linearity makes finding exact, closed-form solutions to the Kushner-Stratonovich equation exceptionally difficult, requiring us to resort to approximations or restrict ourselves to highly specific cases like linear-Gaussian dynamics, which simplifies to the Kalman-Bucy filter.

\section*{Appendix}
\subsection*{A.1. Regularity Regime}
\begin{definition}
For $\alpha\in(0,1)$, the H\"older space $C^\alpha([0,T];\mathbb R)$ consists of continuous functions $z$ such that
\[
[z]_\alpha:=\sup_{0\le s<t\le T}\frac{|z_t-z_s|}{|t-s|^\alpha}<\infty.
\]
\end{definition}

\subsection*{A.2. Rough Integral Obstruction}
\begin{remark}
When $X,Y\in C^\alpha$ with $\alpha\in(1/3,1/2)$ and $Y=f(X)$, the naive Riemann--Stieltjes definition of $\int Y\,dX$ may fail. In the expansion used here, the leading rough term $\mathbb{Z}_{s,t}$ is required.
\end{remark}

\subsection*{A.3. Iterated Integral Object}
\begin{definition}
With our notation,
\[
\mathbb{X}_{s,t}:=X_{s,t}^{(2)}=\int_s^t (X_r-X_s)\,dX_r,
\qquad
\mathbb{Z}_{s,t}:=f'(X_s)\,\mathbb{X}_{s,t}.
\]
So $\mathbb{X}_{s,t}$ denotes the second level of the rough path, and $\mathbb{Z}_{s,t}$ denotes the leading term in the decomposition.
\end{definition}

\subsection*{A.4. Partition and Mesh Size}
\begin{definition}
Formally, a partition of $[s,t]$ is
\[
\mathcal P=\{s=t_0<t_1<\cdots<t_n=t\}.
\]
So an element of $\mathcal P$ is a point $t_i$, not an interval.

The associated subintervals are
\[
[t_{i-1},t_i],\qquad i=1,\dots,n,
\]
and the mesh is
\[
|\mathcal P|
:=
\max_{1\le i\le n}(t_i-t_{i-1}).
\]
\end{definition}

\subsection*{A.5. Young Integration Condition}
\begin{remark}
Young integration is a deterministic pathwise integration theory used to define integrals $\int Y\,dX$, where $Y$ is the integrand and $X$ is the integrator. The paths can be irregular (non-$C^1$) for both $Y$ and $X$, as long as the combined H\"older regularity is high enough.

For Young integration, if $Y\in C^\gamma$ and $X\in C^\alpha$, a sufficient condition for
\[
\int Y\,dX
\]
to be defined is
\[
\alpha+\gamma>1.
\]
In words, $Y$ is H\"older continuous with exponent $\gamma$, and $X$ is H\"older continuous with exponent $\alpha$.
Here $X\in C^\alpha$ with $\alpha<\frac12$, and for $Y=f(X)$ one has the same H\"older order ($\gamma\approx\alpha$). Hence
\[
\alpha+\gamma\approx 2\alpha<1.
\]
So the Young condition fails, and classical pathwise integration may not exist. This is why the rough second-level information (equivalently, the iterated-integral component inside $\mathbb{X}_{s,t}$) is needed.
\end{remark}

\subsection*{A.6. Why $Y=f(X)$ Has the Same H\"older Order}
\begin{remark}
Assume $X\in C^\alpha([0,T])$ and define $Y_t:=f(X_t)$, where $f$ is $C^1$ (in particular, locally Lipschitz on the range of $X$). Since $X$ is continuous on $[0,T]$, its range is compact, so
\[
L:=\sup_{\xi\in X([0,T])}|f'(\xi)|<\infty.
\]
By the mean value theorem, for any $s,t\in[0,T]$,
\[
|Y_t-Y_s|
=
|f(X_t)-f(X_s)|
=
|f'(\xi_{s,t})|\,|X_t-X_s|
\le
L\,|X_t-X_s|,
\]
for some $\xi_{s,t}$ between $X_s$ and $X_t$.
Because $X\in C^\alpha$, there exists $C_X$ such that
\[
|X_t-X_s|\le C_X|t-s|^\alpha.
\]
Hence
\[
|Y_t-Y_s|
\le
LC_X|t-s|^\alpha,
\]
so $Y\in C^\alpha([0,T])$. Therefore $Y=f(X)$ has the same H\"older order as $X$ (equivalently, one can take $\gamma=\alpha$ in the Young condition).
\end{remark}

\subsection*{A.7. Pathwise Meaning of $\int Y\,dX$}
\begin{definition}
Let $X=(X_t)_{t\in[0,T]}$ and $Y=(Y_t)_{t\in[0,T]}$ be stochastic processes on $(\Omega,\mathcal F,\mathbb P)$. A pathwise definition of the integral means: fix $\omega\in\Omega$ first, then define the integral along the deterministic paths $t\mapsto X_t(\omega)$ and $t\mapsto Y_t(\omega)$.

For $s<t$, one writes
\[
\int_s^t Y_r(\omega)\,dX_r(\omega)
:=
\lim_{|\mathcal P|\to0}\sum_{i=0}^{n-1}
Y_{t_i}(\omega)\bigl(X_{t_{i+1}}(\omega)-X_{t_i}(\omega)\bigr),
\]
if the limit exists.

So ``pathwise'' refers to trajectory-by-trajectory construction (for each fixed $\omega$), rather than defining the integral first through probabilistic objects such as expectations or martingale isometries.
After fixing $\omega$, one usually omits $\omega$ from the notation and writes the same formula in compact form as $\int_s^t Y_r\,dX_r$.
\[
\int_s^t Y_r\,dX_r
:=
\lim_{|\mathcal P|\to0}\sum_{i=0}^{n-1}
Y_{t_i}\bigl(X_{t_{i+1}}-X_{t_i}\bigr).
\]
\end{definition}

\subsection*{A.8. Stochastic Process, Path, and Rough Path}
\begin{definition}
A stochastic process is a random map
\[
(t,\omega)\mapsto X_t(\omega).
\]
A path is one realization: fix $\omega$, then
\[
t\mapsto X_t(\omega).
\]

A path is called a rough path when one does not keep only $X$, but an enhanced object
\[
\mathbf X=\bigl(X,\,\mathbb X,\,\mathbf X^{(3)},\dots\bigr),
\]
with enough iterated-integral levels, satisfying algebraic consistency (Chen relation) and analytic regularity bounds.

For $p$-rough paths, one keeps levels up to $\lfloor p\rfloor$. For example, if $2\le p<3$, a rough path is
\[
(X,\mathbb X),
\]
where $\mathbb X$ is the second-level iterated-integral data.
\[
\text{Convention: }\mathbb X_{s,t}=X_{s,t}^{(2)}\text{ denotes the second level of the lift }(X,\mathbb X).
\]
\end{definition}

\subsection*{A.9. Rough Path Levels and Iterated Integrals}
\begin{definition}
For a geometric rough path over
\[
X:[0,T]\to\mathbb R^d,
\]
the levels on $[s,t]$ are:
\[
X_{s,t}^{(0)}=1
\]
(zeroth level, scalar unit),
\[
X_{s,t}^{(1)}=X_t-X_s
\]
(first level, the increment),
\[
X_{s,t}^{(2)}=\int_s^t (X_r-X_s)\otimes dX_r
\]
(second level, iterated integral / L\'evy area-type object),
\[
X_{s,t}^{(3)}
=
\int_s^t\int_s^{r_3}\int_s^{r_2}
dX_{r_1}\otimes dX_{r_2}\otimes dX_{r_3}
\]
(third level, third-order iterated integral).

In general, level $k$ is the $k$-fold iterated integral of $dX$ on $[s,t]$, and rough path theory keeps enough levels up to $\lfloor p\rfloor$ for $p$-rough paths.
For $k\ge2$, these rough path levels are also called iterated-integral levels (or signature levels).
In particular, one often writes the second level as $\mathbb X_{s,t}:=X_{s,t}^{(2)}$.

In words: level $0$ is normalization, level $1$ records net displacement, level $2$ records pairwise interaction/area effects of increments, and level $3$ records third-order interaction effects. Higher levels encode finer path geometry needed to define differential equations driven by rough signals.
For explicit computations, note that since
\[
X_r-X_s=\int_s^r dX_u,
\]
the second level can be written as
\[
X_{s,t}^{(2)}
=
\int_s^t\int_s^r dX_u\otimes dX_r.
\]
For $X=(X^1,\dots,X^d)$, this means componentwise
\[
\bigl(X_{s,t}^{(2)}\bigr)^{i,j}
=
\int_s^t (X_r^i-X_s^i)\,dX_r^j,
\qquad i,j=1,\dots,d.
\]
For the third level,
\[
X_{s,t}^{(3)}
=
\int_s^t\int_s^{r_3}\int_s^{r_2}
dX_{r_1}\otimes dX_{r_2}\otimes dX_{r_3},
\]
the componentwise form is
\[
\bigl(X_{s,t}^{(3)}\bigr)^{i,j,k}
=
\int_s^t\int_s^{r_3}\int_s^{r_2}
dX_{r_1}^i\,dX_{r_2}^j\,dX_{r_3}^k,
\qquad i,j,k=1,\dots,d.
\]
So level $2$ takes values in $\mathbb R^d\otimes\mathbb R^d$ (matrix-type object), and higher levels take values in higher-order tensor spaces.
\end{definition}

\subsection*{A.10. Outer (Kronecker) Product}
\begin{definition}
The symbol $\otimes$ denotes the outer (Kronecker) product. If $a,b\in\mathbb R^d$, then
\[
a\otimes b\in\mathbb R^{d\times d},
\qquad
(a\otimes b)_{ij}=a_i b_j.
\]
\end{definition}

\subsection*{A.11. Regularity Comparison: Leading Term vs. Remainder}
\begin{lemma}
Assume $\alpha\in(1/3,1/2)$, $X\in C^\alpha([0,T];\mathbb R^d)$, and $f\in C^2$ with bounded $f''$ on the range of $X$. Define
\[
\mathbb Z_{s,t}:=f'(X_s)\int_s^t (X_r-X_s)\,dX_r,
\qquad
\widetilde R_{s,t}
:=
\int_s^t\int_s^r \bigl(f'(X_u)-f'(X_s)\bigr)\,dX_u\,dX_r.
\]
Then there exists $C>0$ (independent of $s,t$) such that
\[
|\mathbb Z_{s,t}|\le C|t-s|^{2\alpha},
\qquad
|\widetilde R_{s,t}|\le C|t-s|^{3\alpha}.
\]
In particular, the remainder is more regular than the leading rough term.
\end{lemma}
\begin{proof}
Set
\[
\|X\|_\alpha:=[X]_{C^\alpha([0,T])}
:=
\sup_{u\ne v}\frac{|X_u-X_v|}{|u-v|^\alpha}.
\]
Let $(X,\mathbb X)$ be a fixed rough-path lift of regularity $\alpha$.

For the leading term, define $\Phi_r:=X_r-X_s$ on $[s,t]$. Then $\Phi\in C^\alpha$ and
\[
\,[\Phi]_{C^\alpha([s,t])}\le \|X\|_\alpha.
\]
Indeed, for $r,q\in[s,t]$,
\[
\frac{|\Phi_r-\Phi_q|}{|r-q|^\alpha}
=
\frac{|(X_r-X_s)-(X_q-X_s)|}{|r-q|^\alpha}
=
\frac{|X_r-X_q|}{|r-q|^\alpha}
\le
\|X\|_\alpha.
\]
Since $\alpha<\frac12$, one has $2\alpha<1$, so Young integration does not apply here. In the rough-path setting, one interprets the integral with a lift $(X,\mathbb X)$ and uses the rough integration estimate (for controlled integrands), which yields
\[
\left|\int_s^t \Phi_r\,dX_r\right|
\le
C\bigl(\|X\|_\alpha^2+\|\mathbb X\|_{2\alpha}\bigr)|t-s|^{2\alpha}
\le
C|t-s|^{2\alpha}.
\]
Since $X([0,T])$ is compact and $f'$ is continuous, $\sup_{\xi\in X([0,T])}|f'(\xi)|<\infty$.
Hence
\[
|\mathbb Z_{s,t}|
\le
|f'(X_s)|\,C\bigl(\|X\|_\alpha^2+\|\mathbb X\|_{2\alpha}\bigr)|t-s|^{2\alpha}
\le
C|t-s|^{2\alpha}.
\]

For the remainder, set $H_u:=f'(X_u)-f'(X_s)$. Since $f''$ is bounded on $X([0,T])$,
\[
|H_u-H_v|
\le
\|f''\|_\infty |X_u-X_v|
\le
\|f''\|_\infty \|X\|_\alpha |u-v|^\alpha,
\]
so $H\in C^\alpha([s,t])$. Also,
\[
|H_u|
\le
\|f''\|_\infty |X_u-X_s|
\le
\|f''\|_\infty \|X\|_\alpha |u-s|^\alpha.
\]
Define
\[
G_r:=\int_s^r H_u\,dX_u.
\]
This is again a rough integral (since $\alpha+\alpha<1$). The rough integration estimate gives
\[
|G_r-G_q|
\le
C|r-q|^{2\alpha},
\]
so $G\in C^{(2\alpha)}([s,t])$, and in particular
\[
|G_r|=|G_r-G_s|\le C|r-s|^{2\alpha}.
\]
Now
\[
\widetilde R_{s,t}=\int_s^t G_r\,dX_r.
\]
Since $G\in C^{(2\alpha)}$ and $X\in C^\alpha$ with $2\alpha+\alpha=3\alpha>1$, Young's estimate gives
\[
|\widetilde R_{s,t}|
\le
C\,[G]_{C^{(2\alpha)}([s,t])}\,[X]_{C^\alpha([s,t])}|t-s|^{3\alpha}
\le
C|t-s|^{3\alpha}.
\]
Therefore the remainder has higher order than the leading term: $3\alpha>2\alpha$.
\end{proof}

\subsection*{A.12. Taylor-Type Interpretation of Rough Integration}
\begin{remark}
Rough path integration is fundamentally a Taylor-type expansion of increments.

Locally, one expands
\[
\int_s^t f(X_r)\,dX_r
=
\underbrace{f(X_s)(X_t-X_s)}_{\text{first order }|t-s|^{\alpha}}
+\underbrace{f'(X_s)\mathbb{X}_{s,t}}_{\text{second order }|t-s|^{2\alpha}}
+\underbrace{\widetilde R_{s,t}}_{\text{higher order }|t-s|^{3\alpha}}.
\]
Here:
\begin{itemize}
\item first term = first-order part;
\item second term = second-order correction (iterated integral);
\item remainder = higher-order term (here of order $\sim |t-s|^{3\alpha}$).
\end{itemize}

So the method is very similar in spirit to Taylor expansion, but adapted to rough signals via iterated integrals and Chen's algebraic structure.
\end{remark}

\subsection*{A.13. Proof of the Remainder Bound Used in Section 5.2}
\begin{lemma}
Assume $\alpha\in(1/3,1/2)$, $X\in C^\alpha([0,T])$, and a rough-path lift $(X,\mathbb X)$ with $\mathbb X\in C^{(2\alpha)}$. Let $f\in C_b^2(\mathbb R)$ and define
\[
\widetilde R_{s,t}
:=
\int_s^t f(X_r)\,dX_r
-f(X_s)(X_t-X_s)
-f'(X_s)\mathbb X_{s,t}.
\]
Then there exists $K>0$ such that
\[
|\widetilde R_{s,t}|
\le
K\|f\|_{C_b^2}
\Bigl(
\|X\|_{C^\alpha}^3
+\|X\|_{C^\alpha}\|\mathbb X\|_{C^{(2\alpha)}}
\Bigr)|t-s|^{3\alpha}.
\]
\end{lemma}
\begin{proof}
Set
\[
\rho_{s,r}:=f(X_r)-f(X_s)-f'(X_s)(X_r-X_s).
\]
Because $\rho_{s,r}$ is exactly the second-order Taylor remainder, Taylor's theorem (mean-value remainder form) gives, for some $\xi$ between $X_s$ and $X_r$,
\[
\rho_{s,r}
=
\frac12 f''(\xi)\,(X_r-X_s)^2.
\]
Hence
\[
|\rho_{s,r}|
\le
\frac12\sup_x |f''(x)|\,|X_r-X_s|^2
\le
\|f\|_{C_b^2}|X_r-X_s|^2
\le
\|f\|_{C_b^2}\|X\|_{C^\alpha}^2|r-s|^{2\alpha}.
\]
Hence
\[
[\rho]_{C^{(2\alpha)}([s,t])}
\le
C\|f\|_{C_b^2}\|X\|_{C^\alpha}^2.
\]
More explicitly, for $u<v$ in $[s,t]$,
\[
\frac{|\rho_{u,v}|}{|v-u|^{2\alpha}}
\le
\frac{\|f\|_{C_b^2}|X_v-X_u|^2}{|v-u|^{2\alpha}}
\le
\|f\|_{C_b^2}\|X\|_{C^\alpha}^2,
\]
and taking the supremum over $u<v$ gives the stated bound.
Also
\[
[f(X)]_{C^\alpha([s,t])}
\le
\|f\|_{C_b^1}\|X\|_{C^\alpha}.
\]
Indeed, for $u<v$ in $[s,t]$,
\[
\frac{|f(X_v)-f(X_u)|}{|v-u|^\alpha}
\le
\|f'\|_{L^\infty}\frac{|X_v-X_u|}{|v-u|^\alpha}
\le
\|f\|_{C_b^1}\|X\|_{C^\alpha},
\]
and taking the supremum over $u<v$ yields the claim.

Write the increment decomposition
\[
f(X_r)-f(X_s)=f'(X_s)X_{s,r}+\rho_{s,r}.
\]
Applying the rough integration (sewing) estimate to this controlled structure gives
\[
|\widetilde R_{s,t}|
\le
C\Bigl(
[\rho]_{C^{(2\alpha)}([s,t])}[X]_{C^\alpha([s,t])}
+[f(X)]_{C^\alpha([s,t])}[\mathbb X]_{C^{(2\alpha)}([s,t])}
\Bigr)|t-s|^{3\alpha}.
\]
Substituting the two bounds above and absorbing constants,
\[
|\widetilde R_{s,t}|
\le
K\|f\|_{C_b^2}
\Bigl(
\|X\|_{C^\alpha}^3
+\|X\|_{C^\alpha}\|\mathbb X\|_{C^{(2\alpha)}}
\Bigr)|t-s|^{3\alpha}.
\]
\end{proof}

\subsection*{A.14. Norm Conventions Used in the Remainder Estimate}
\begin{definition}
The $C_b^2$-norm of $f$ is
\[
\|f\|_{C_b^2}
:=
\sup_x |f(x)|
+\sup_x |f'(x)|
+\sup_x |f''(x)|.
\]

The $\alpha$-H\"older seminorm of $X$ is
\[
[X]_{C^\alpha([0,T])}
:=
\sup_{0\le s<t\le T}\frac{|X_t-X_s|}{|t-s|^\alpha}.
\]

The $\alpha$-H\"older norm of the path $X$ is
\[
\|X\|_{C^\alpha([0,T])}
:=
\sup_{t\in[0,T]}|X_t|
+\sup_{0\le s<t\le T}\frac{|X_t-X_s|}{|t-s|^\alpha}.
\]
Equivalently,
\[
\|X\|_{C^\alpha([0,T])}
=
\sup_{t\in[0,T]}|X_t|+[X]_{C^\alpha([0,T])}.
\]

The $2\alpha$-H\"older seminorm of $\mathbb X$ is
\[
[\mathbb X]_{C^{(2\alpha)}([0,T]^2)}
:=
\sup_{0\le s<t\le T}\frac{|\mathbb X_{s,t}|}{|t-s|^{2\alpha}}.
\]

The $2\alpha$-H\"older norm of the second-level increment map $\mathbb X$ is
\[
\|\mathbb X\|_{C^{(2\alpha)}([0,T]^2)}
:=
\sup_{0\le s<t\le T}\frac{|\mathbb X_{s,t}|}{|t-s|^{2\alpha}}.
\]
In this note, for $\mathbb X$ we use this seminorm-type quantity.
\end{definition}

\subsection*{A.15. It\^o's Formula}
\begin{theorem}[It\^o's formula, one-dimensional]
Let $X_t$ satisfy
\[
dX_t=b_t\,dt+\sigma_t\,dW_t,
\]
and let $\varphi\in C^{1,2}([0,T]\times\mathbb R)$. Then
\[
d\varphi(t,X_t)
=
\partial_t\varphi(t,X_t)\,dt
+\partial_x\varphi(t,X_t)\,dX_t
+\frac12\,\partial_{xx}\varphi(t,X_t)\,\sigma_t^2\,dt.
\]
Equivalently, for $0\le s\le t\le T$,
\begin{align*}
\varphi(t,X_t)-\varphi(s,X_s)
&=
\int_s^t \partial_t\varphi(r,X_r)\,dr
+\int_s^t \partial_x\varphi(r,X_r)\,dX_r \\
&\quad
+\frac12\int_s^t \partial_{xx}\varphi(r,X_r)\,\sigma_r^2\,dr.
\end{align*}
\end{theorem}

\begin{remark}[Brownian case used in Section 5.3]
For Brownian motion $W$ and $\varphi(x)=\frac12 x^2$, It\^o's formula gives
\[
d\!\left(\frac12 W_t^2\right)=W_t\,dW_t+\frac12\,dt.
\]
Integrate from $s$ to $t$:
\[
\frac12\bigl(W_t^2-W_s^2\bigr)
=
\int_s^t W_r\,dW_r+\frac12(t-s).
\]
So
\[
\int_s^t W_r\,dW_r
=
\frac12\bigl(W_t^2-W_s^2\bigr)-\frac12(t-s).
\]
The term $-\frac12(t-s)$ comes from Brownian quadratic variation $[W]_t=t$.
This is the identity used in the It\^o computation of
\[
\int_s^t (X_r-X_s)\,dX_r.
\]
\end{remark}

\subsection*{A.16. Algebraic Consistency and Analytic Regularity}
\begin{definition}
For a two-level rough path $(X,\mathbb X)$ with $\alpha\in(1/3,1/2)$, the two structural requirements are:

\textbf{Algebraic consistency (Chen relation).} For all $s\le u\le t$,
\[
\mathbb X_{s,t}
=
\mathbb X_{s,u}
+\mathbb X_{u,t}
+X_{s,u}\otimes X_{u,t},
\qquad
X_{s,u}:=X_u-X_s.
\]
In words, the second-level increment on $[s,t]$ must decompose consistently when the interval is split at an intermediate point $u$.

\textbf{Analytic regularity bounds.} There exists $C<\infty$ such that for all $s<t$,
\[
|X_{s,t}|\le C|t-s|^\alpha,
\qquad
|\mathbb X_{s,t}|\le C|t-s|^{2\alpha}.
\]
Equivalently, in seminorm form,
\[
[X]_{C^\alpha}<\infty,
\qquad
[\mathbb X]_{C^{(2\alpha)}}<\infty.
\]
In words, the first level scales like an $\alpha$-H\"older increment and the second level scales like a $2\alpha$-H\"older increment.
These two conditions together are what make rough Riemann sums stable under partition refinement.
\end{definition}

\subsection*{A.17. Gaussian Random Field}
\begin{definition}
A Gaussian random field is a stochastic process indexed by space (or another index set) such that every finite vector of values is jointly Gaussian.

Formally, for a field $(X(x))_{x\in D}$ on a domain $D$, one requires that for any $n\in\mathbb N$ and any $x_1,\dots,x_n\in D$,
\[
\bigl(X(x_1),\dots,X(x_n)\bigr)
\]
has a multivariate normal distribution.

It is characterized by:
\[
m(x):=\mathbb E[X(x)],
\qquad
K(x,y):=\mathrm{Cov}(X(x),X(y)).
\]
In the stationary linear SPDE example above, the invariant law at fixed time is a mean-zero Gaussian random field with covariance operator $\frac{\sigma^2}{2}(I-\partial_{xx})^{-1}$.
\end{definition}

\subsection*{A.18. Fourier-Series Representation of the Covariance Kernel}
\begin{remark}
Let
\[
C:=\frac{\sigma^2}{2}(I-\partial_{xx})^{-1},
\qquad
(Cf)(x)=\int_0^{2\pi}K(x-y)f(y)\,dy.
\]
Equivalently, in convolution notation on the torus, $Cf=K*f$.
On the periodic domain, Fourier modes diagonalize $C$, so
\[
\widehat K(k)=\frac{\sigma^2}{2}\frac{1}{1+k^2}\cdot\frac{1}{2\pi}.
\]
By inverse Fourier series,
\[
K(x)=\sum_{k\in\mathbb Z}\widehat K(k)e^{ikx}
=\frac{\sigma^2}{4\pi}\sum_{k\in\mathbb Z}\frac{e^{ikx}}{1+k^2}.
\]
Using $\widehat K(-k)=\widehat K(k)$,
\[
K(x)=\frac{\sigma^2}{2\pi}\sum_{k\in\mathbb Z}\frac{\cos(kx)}{1+k^2}.
\]
\end{remark}

\subsection*{A.19. What the $k$-th Fourier Mode Means}
\begin{definition}
The $k$-th Fourier mode means the basis oscillation with frequency $k$:
\[
\text{complex form: } e^{ikx},
\qquad
\text{real form: } \cos(kx),\ \sin(kx).
\]
So ``$k$-th mode'' is the component of a function at wave number $k$.

In Fourier series,
\[
f(x)=\sum_{k\in\mathbb Z}\widehat f(k)e^{ikx},
\]
and
\[
\widehat f(k)e^{ikx}
\]
is the $k$-th mode contribution.
\end{definition}

\subsection*{A.20. Why Use the Kernel Representation}
\begin{remark}
The covariance operator form
\[
C=\frac{\sigma^2}{2}(I-\partial_{xx})^{-1}
\]
is abstract. Writing
\[
(Cf)(x)=\int_0^{2\pi}K(x-y)f(y)\,dy
\]
makes the spatial structure explicit. In particular,
\[
\mathrm{Cov}(X(x),X(y))=K(x-y)
\]
for the stationary mean-zero field. So the kernel representation is what directly tells us how fluctuations at points $x$ and $y$ are correlated in space, and it is the form used to derive closed formulas from Fourier coefficients.
\end{remark}

\end{document}
